\section{Conclusion}

This work investigated the use of a hierarchical Bayesian model (HBM) for predicting battery end-of-life (EoL) from early-cycle features, incorporating cycling condition as a group-level structure rather than a flat input feature. Three early-cycle features were extracted following the approach of \cite{severson2019}, and cells were partitioned into 8 groups based on charging protocol using a hill-climbing algorithm, which produced a variance partition coefficient sufficient to justify a hierarchical treatment of the data.

Across 20 cross-validation folds, the HBM achieved a mean MAPE of $14.7\%$, outperforming both the plain ridge baseline ($22.5\%$) and the ridge model with the group feature ($16.8\%$). On the held-out test set, the HBM maintained the lowest MAPE at $13.9\%$, though its RMSE was marginally higher than the baselines, suggesting sensitivity to a small number of outlier predictions. MCMC diagnostics confirmed reliable posterior inference, with all parameters satisfying $\hat{R} \leq 1.01$ and adequate effective sample sizes.

Bayesian posterior analysis revealed that the hyperprior slope linking mean C-rate to the group intercept was credibly negative (HDI $[-4.91, -1.55]$), providing probabilistic evidence that higher charge rates are associated with reduced battery lifetime at the group level. The early-cycle feature coefficients, by contrast, were found to generalize across groups, with no credible C-rate modulation. Together, these findings suggest that the performance advantage of the HBM over conventional regression stems not from learning different feature sensitivities per group, but from appropriately pooling information across cells sharing similar cycling conditions when estimating baseline lifetime.

Future work could explore richer group structures and the influence of cluster size and number. The grouping of the cycling conditions heavily influenced the results of HBM as well as the baseline rigde regression with the group feature.